\documentclass[11pt]{article}
\usepackage[margin=0.92in]{geometry}
\usepackage{amsmath,amssymb,amsthm,booktabs,microtype}
\usepackage[colorlinks=true,linkcolor=black,urlcolor=blue]{hyperref}
\newtheorem{theorem}{Theorem}
\newtheorem{lemma}{Lemma}
\newcommand{\Q}{\mathbb Q}
\newcommand{\Z}{\mathbb Z}
\newcommand{\vp}{v_p}
\title{\bfseries Seven terms for ten singular moduli}
\author{An exact-arithmetic supplement to the ProveIt transfer reports}
\date{26 November 2026}
\begin{document}
\maketitle
\noindent
For a singular modulus $j$ of discriminant $\Delta$, write $s(\Delta)$ for
the smallest number of nonzero terms of a nonzero $P\in\Q[X]$ with $P(j)=0$.
The number is independent of the chosen conjugate. The exponents in $P$
are unrestricted nonnegative integers.

\begin{theorem}
$s(\Delta)=7$ for
\[
 \Delta\in\{-152,-244,-300,-451,-472,-588,-628,-867,-1083,-1147\}.
\]
\end{theorem}
\noindent
In particular, this proves the result for the original discriminant-$-152$
modulus, not just its powers divisible in exponent by $182$ as in the
\href{https://chatgpt.com/c/6b07cfad-a4b4-83e8-add3-8b154f1a8176?messageId=77e95482-c17b-54ea-bd2e-3033fb0d7dfa}{transfer report}.
No Riemann hypothesis or bound on the largest exponent is used. The result
does not settle the original modulus of discriminant $-87$.

\section*{A local six-term obstruction}
Normalize $\vp(p)=1$ and work in a finite extension of $\Q_p$ containing all
roots under discussion, with valuation ring $\mathcal O$. For a nonzero center
$c$ and an integer $k>1/(p-1)$, the map $Y\mapsto c\exp(p^kY)$ is a bijection
from $\mathcal O$ to $c(1+p^k\mathcal O)$. Its inverse uses the $p$-adic
logarithm. Set $w(n)=kn-\vp(n!)$. If
\[
             w(n)>\max_{0\le i\le5}w(i) \qquad(n\ge6),
             \tag{1}
\]
every polynomial with at most six terms has at most five distinct roots
in that disk.

Indeed, for $P(X)=\sum_{j=1}^m a_jX^{e_j}$, $1\le m\le6$, put
$A_n=\sum_j a_jc^{e_j}e_j^n$, with $e_j^0=1$. The Vandermonde determinant
shows that $M=\max_{i<m}|A_i|_p$ is positive. The monic polynomial
$\prod_j(T-e_j)$ gives a recurrence for the $A_n$ with integral coefficients,
so $|A_n|_p\le M$ for every $n$. The coefficient of $Y^n$ in
$P(c\exp(p^kY))$ has norm $|A_n|_p p^{-w(n)}$. At an index $i<m$ attaining
$M$, it has norm $M p^{-w(i)}$; by (1), every coefficient with index
at least six has smaller norm. The coefficients tend to zero, since
$k>1/(p-1)$. Strassmann's theorem bounds the number of roots by the
last index of maximum coefficient norm, hence by five.

\section*{Ten integral disk certificates}
Let $H_\Delta$ denote the degree-six ring class polynomial of $\Delta$.
Its six roots are distinct and conjugate over $\Q$ by the standard
complex multiplication theorem. Write
\[
 H_\Delta(p^a(u+Y))=\sum_{i=0}^6 B_iY^i .
\]
Each row below gives $\vp(B_i)\ge 6a+k(6-i)$ and $B_6=p^{6a}$, with
$p\nmid u$. Therefore $p^{-6(a+k)}H_\Delta(p^a(u+p^kY))$ is monic
with integral coefficients. All its roots are integral: if $\vp(Y)<0$,
the leading term has strictly smaller valuation than every lower term.
It follows that every root of $H_\Delta$ lies in
$p^au(1+p^k\mathcal O)$.

\medskip
\begin{center}
\begin{tabular}{@{}rrrrrl@{}}
\toprule
$\Delta$&$p$&$a$&$u$&$k$&\multicolumn{1}{l}{$\vp(B_0),\ldots,\vp(B_6)$}\\
\midrule
$-152$&2&6&1&2&$48,47,44,44,40,39,36$\\
$-244$&2&6&3&2&$48,49,44,46,40,41,36$\\
$-300$&2&4&3&2&$36,35,32,32,28,27,24$\\
$-451$&3&3&1&2&$30,30,27,24,26,21,18$\\
$-472$&2&6&1&2&$48,47,44,44,40,39,36$\\
$-588$&2&4&3&2&$36,35,32,33,28,27,24$\\
$-628$&2&6&3&2&$48,49,44,46,40,41,36$\\
$-867$&3&3&2&2&$32,30,27,25,22,23,18$\\
$-1083$&17&0&8&1&$6,5,4,4,2,1,0$\\
$-1147$&3&3&1&2&$30,29,27,24,23,21,18$\\
\bottomrule
\end{tabular}
\end{center}

For $(p,k)=(2,2)$, the maximum through index five is $7$;
$w(6)=8$ and $w(n)=n+s_2(n)\ge8$ for $n\ge7$, with $s_p(n)$ the
base-$p$ digit sum. For $(3,2)$, the maximum is $9$ and
$w(n)=(3n+s_3(n))/2>9$ for $n\ge6$. For $(17,1)$, the maximum is $5$ and
$w(n)\ge15n/16>5$ for $n\ge6$. Thus each row satisfies (1).
A rational annihilator of one conjugate vanishes at all six, so cannot
have at most six terms. The class polynomial itself gives at most seven.
This proves the theorem.

\paragraph{Reproduction and provenance.}
The coefficient lists used here come from
\href{https://github.com/VladimirReshetnikov/ProveIt/blob/48efe71a726dddc5a505534edaa4c0fd190416b9/docs/incoming/grh_sparse_singular_moduli.zip}{the pinned ProveIt archive},
\texttt{checks/five\_small\_witnesses.json}. Its \texttt{local.tex} supplies
the local argument, related to
\href{https://arxiv.org/abs/1008.4808}{Avenda\~no--Krick, \emph{Sharp bounds for the number of roots of univariate fewnomials}}.
The accompanying checker recomputes the table by binomial expansion and
cross-checks it by modular Horner substitution. It verifies the
divisibility certificates from the archived class polynomials; it does
not recompute the class polynomials from complex multiplication.
\end{document}

